\begin{textcolorbox}[Prompt for Rephrasing Proofs to Relation Problems]
\textbf{Task:} Transform the given inequality proof problem into a relation prediction problem.
\\\\
\textbf{Instructions:}

1. Analyze the original problem, identifying key components such as variables, domains, conditions, and the main inequality. \\
2. Rephrase the problem by maintaining the original expressions and replacing the relation symbol with a blank to be filled. \\
3. Preserve any additional conditions or constraints from the original problem in your rephrased version. \\
4. Change the task from ``Prove'' to ``Determine the correct inequality relation to fill in the blank.'' \\
5. Provide a set of options for the relation, always including $\leq$, $\geq$, $=$, $<$, $>$, and ``None of the above''. \\
6. Determine the correct answer based on your modification and analysis. \\

\textbf{Output Format:}

Provide your response in the following structure:

\vspace{2mm}
\texttt{<Analysis>}: Detailed step-by-step analysis of the original problem and your approach to rephrasing it.

\vspace{2mm}
\texttt{<Conclusion>}: YES or NO, followed by a brief explanation of whether and how the problem can be effectively rephrased.

\vspace{2mm}
\texttt{<Rephrased Problem>}: 

\quad Transformed problem statement.

\quad Options:
(A) $\leq$ \quad
(B) $\geq$\quad
(C) $=$ \quad
(D) $<$ \quad
(E) $>$ \quad
(F) None of the above

\vspace{2mm}
\texttt{<Answer>}: Option letter.
\\\\
\textbf{Key Considerations:} \\
1. Maintain the original mathematical expressions and any given conditions as much as possible.\\
2. Ensure the rephrased problem captures the essence and complexity of the original problem.\\
3. For problems with multiple inequalities, focus on one main inequality for the relation prediction task.\\
4. When dealing with complex fractions or expressions, keep them intact to maintain the problem's difficulty level.\\
5. If the relation depends on specific values of the variables or cannot be definitively determined, consider using ``None of the above'' as the correct answer.
% 6. Use LaTeX format ($$) for all variables and formulas, not \\( and \\).
% ations, use double dollar signs ($$) on separate lines.
\\\\
\textbf{Example:} \\
% \texttt{\{examples\}}
Original problem:
Let $a, b, c \in \mathbb{R}^{+}$. Prove the inequality
$$
\frac{a b c}{(1+a)(a+b)(b+c)(c+16)} \leq \frac{1}{81}
$$
\vspace{2mm}
\texttt{<Analysis>}: To rephrase it to a relation prediction problem, we can focus on the following steps: \\
1. The original problem is a proof task for an inequality involving positive real numbers $a$, $b$, and $c$. \\
2. The left side of the inequality is a complex fraction $\frac{a b c}{(1+a)(a+b)(b+c)(c+16)}$. \\
 % with a product in the numerator and a product of sums in the denominator
3. The right side is a constant fraction $\frac{1}{81}$. \\
4. The original inequality uses the ``less than or equal to'' ($\leq$) relation, which needs to hold for all positive real values of $a$, $b$, and $c$. \\
5. We can transform the proof task into determining the correct relation between the left and right sides of the inequality.

\vspace{2mm}
\texttt{<Conclusion>}: YES, the inequality can be effectively rephrased as a relation prediction problem.

\vspace{2mm}
\texttt{<Rephrased problem>}: \\
Let $a, b, c \in \mathbb{R}^{+}$. Consider the following inequality:
$$
\frac{a b c}{(1+a)(a+b)(b+c)(c+16)} \quad (\quad) \quad \frac{1}{81} .
$$
Determine the correct inequality relation to fill in the blank. 
\\
Options:
(A) $\leq$ \quad
(B) $\geq$\quad
(C) $=$ \quad
(D) $<$ \quad
(E) $>$ \quad
(F) None of the above

\vspace{2mm}
\texttt{<Answer>:} A
\\\\
\textbf{Now, please rewrite the following problem:}

Original problem: \texttt{\{problem\}}
\end{textcolorbox}